\subsection{点分布与切空间}

在本节中，X 表示一个 C \(^{r}\) 类流形。

13.3.1. 设 \(x \in X\)，并设 \(k\) 是一个 \(\leqslant r\) 的整数。设 \(c = (U, \varphi, E)\) 是 \(X\) 的以 \(x\) 为中心的一个卡；13.2.1 中定义的同构 \(\theta_c \colon \mathsf{TS}^{(r)}(E) \to \mathsf{T}_x^{(r)}(X)\) 把 \(\mathsf{TS}^{(k)}(E)\) 映到 \(\mathsf{T}_x^{(k)}(X)\) 上，把 \(\mathsf{TS}^{(k-1)}(E)\) 映到 \(\mathsf{T}_x^{(k-1)}(X)\) 上；通过取限制和取商，它诱导出同构

\[
\theta_ {c, k} \colon \mathsf {T S} ^ {k} (\mathrm{E}) = \mathsf {T S} ^ {(k)} (\mathrm{E}) / \mathsf {T S} ^ {(k - 1)} (\mathrm{E}) \rightarrow \mathsf {T} _ {x} ^ {(k)} (\mathrm{X}) / \mathsf {T} _ {x} ^ {(k - 1)} (\mathrm{X}).
\]

另一方面，c 定义了一个同构，仍记作 \(\theta_{c}\)（参见 5.5.1），它把 E 映到切空间 \(T_{x}(X)\) 上。用 \(i_{k}\) 表示复合映射

\[
\mathrm{T} _ {x} ^ {(k)} (\mathrm{X}) / \mathrm{T} _ {x} ^ {(k - 1)} (\mathrm{X}) \xrightarrow {\theta_ {c , k} ^ {- 1}} \mathrm{TS} ^ {k} (\mathrm{E}) \xrightarrow {\mathrm{TS} ^ {k} (\theta_ {c})} \mathrm{TS} ^ {k} (\mathrm{T} _ {x} (\mathrm{X})).
\]

同构 \(i_k\) 与 \(c\) 的选取无关；当 \(k = 1\) 时，用它把 \(\mathrm{T}_x^{(1)+}(\mathrm{X}) = \mathrm{T}_x^{(1)}(\mathrm{X}) / \mathrm{T}_x^{(0)}(\mathrm{X})\) 与切空间 \(\mathrm{T}_x(\mathrm{X})\) 恒同。若 \(t \in \mathrm{T}_x^{(k)}(\mathrm{X})\)，可以把 \(i_k(t)\) 记为 \(i_k\) 作用于 \(t\) 模 \(\mathrm{T}_x^{(k-1)}(\mathrm{X})\) 的类而得的像。

令

\[
\operatorname{gr} _ {k} \mathrm{T} _ {x} ^ {(r)} (\mathrm{X}) = \mathrm{T} _ {x} ^ {(k)} (\mathrm{X}) / \mathrm{T} _ {x} ^ {(k - 1)} (\mathrm{X}) \quad \text { and } \quad \operatorname{gr} \mathrm{T} _ {x} ^ {(r)} (\mathrm{X}) = \bigoplus_ {0 \leqslant k \leqslant r} \operatorname{gr} _ {k} \mathrm{T} ^ {(r)} (\mathrm{X}).
\]

我们称 gr \(\mathbf{T}_{x}^{(r)}(\mathbf{X})\) 为与滤链 \((\mathbf{T}_{x}^{(k)}(\mathbf{X}))_{k\leq r}\)（\(\mathbf{T}_{x}^{(r)}(\mathbf{X})\) 的递增滤链）相伴的分次对象。诸 \(i_k\) 定义了一个分次向量空间的同构

\[
i \colon \operatorname{gr} \mathrm{T} _ {x} ^ {(r)} (\mathrm{X}) \rightarrow \mathrm{TS} ^ {(r)} (\mathrm{T} _ {x} (\mathrm{X})).
\]

若 \(r = \infty\) 或 \(\omega\)，则 \(i\) 是 gr \(T_x^{(r)}(X)\) 到 \(TS(T_x(X))\) 上的同构。当希望指明 \(x\)（或 X，或两者）时，就写 \(i_x\)（或 \(i_X\)，或 \(i_{X,x}\)）以代替 \(i\)。

13.3.2. 在上述假设之外，再假设 X 局部有限维。相对于 X 各点的诸同构 \(i_{k}\) 定义了一个向量丛同构

\[
i _ {k} \colon \mathrm{T} ^ {(k)} (\mathrm{X}) / \mathrm{T} ^ {(k - 1)} (\mathrm{X}) \rightarrow \mathrm{TS} ^ {k} (\mathrm{T} (\mathrm{X}))
\]

其中 TS \(^{k}\) (T(X)) 表示 C \(^{r-1}\) 类的、由有限维向量函子 TS \(^{k}\) 从 T(X) 得到的向量丛（参见 7.6.5）。当 \(k \geqslant 1\) 时，\(i_{k}\) 是一个 C \(^{r-k}\) 类同构；当 \(k = 0\) 时，它是平凡丛 \(\mathrm{K}_{\mathrm{X}}\) 的恒同映射。

13.3.3. 例。—— 在上述假设下，设 \(\xi = (\xi^{1}, \ldots, \xi^{n})\) 是 X 在 \(x\) 处的一个坐标系。用 \((\partial_{i,x})_{1 \leqslant i \leqslant n}\) 表示 \(\mathrm{T}_{x}(\mathrm{X})\) 的由 \(\xi\) 定义的基（参见 5.5.8），并用 \(((\Delta_{\xi}^{\infty})_{x})_{|\alpha| \leqslant k}\) 表示 13.2.6 中定义的 \(\mathrm{T}_{x}^{(k)}(\mathrm{X})\) 的基。我们有：

\[
\begin{array}{l l} i _ {k} ((\Delta_ {\xi} ^ {\alpha}) _ {x}) = 0 & \text {if} \quad | \alpha | <   k \\ i _ {k} ((\Delta_ {\xi} ^ {\alpha}) _ {x}) = \gamma_ {\alpha_ {1}} (\partial_ {1, x}) \dots \gamma_ {\alpha_ {n}} (\partial_ {n, x}) & \text {if} \quad | \alpha | = k. ^ {(1)} \end{array}
\]

\(^{1}\) 这里同样，诸 \(\gamma_{\alpha_{i}}(\partial_{i,x})\) 的乘积是 \(\mathrm{TS}(\mathrm{T}_{x}(\mathrm{X}))\) 中的对称积。

13.3.4. 假设 K 的特征为零，或者 X 局部有限维。设 k 是满足 \(0 \leqslant k \leqslant r\) 的整数，并设 \(x \in X\)；设 F 是一个 Banach 空间。根据 12.6.8，我们有正合序列

\[
0 \rightarrow \mathrm{P} _ {k} (\mathrm{T} _ {x} (\mathrm{X}); \mathrm{F}) \stackrel {{\iota}} {{\rightarrow}} \mathrm{P} _ {x} ^ {k} (\mathrm{F} _ {\mathrm{X}}) \stackrel {{\sigma}} {{\rightarrow}} \mathrm{P} _ {x} ^ {k - 1} (\mathrm{F} _ {\mathrm{X}}) \rightarrow 0,\tag{i}
\]

其中 \(\sigma = r^{k, k-1}\)。另一方面，根据 13.3.1，我们有正合序列

\[
0 \leftarrow \mathsf {T S} ^ {k} (\mathrm{T} _ {x} (\mathrm{X})) \stackrel {{i}} {{\leftarrow}} \mathrm{T} _ {x} ^ {(k)} (\mathrm{X}) \stackrel {{s}} {{\leftarrow}} \mathrm{T} _ {x} ^ {(k - 1)} (\mathrm{X}) \leftarrow 0,\tag{ii}
\]

其中 s 是 \(\mathrm{T}_{x}^{(k-1)}(\mathrm{X})\) 到 \(\mathrm{T}_{x}^{(k)}(\mathrm{X})\) 的包含映射，且 \(i=i_{k}\)。这些正合序列在 F 中“成对”。更确切地说，把它们写成如下形式：

\begin{enumerate}
\def\labelenumi{(\roman{enumi})}
\tightlist
\item
\end{enumerate}

\[
0 \rightarrow \mathrm{A} \xrightarrow {\iota} \mathrm{B} \xrightarrow {\sigma} \mathrm{C} \rightarrow 0\tag{ii}
\]

\[
0 \leftarrow A ^ {\prime} \stackrel {i} {\leftarrow} B ^ {\prime} \stackrel {s} {\leftarrow} C ^ {\prime} \leftarrow 0.
\]

若 \(a \in \mathbf{A}\) 且 \(a' \in \mathbf{A}'\)，F 中的元素 \(\langle a, a' \rangle\) 由 13.1.2 定义；若 \(b \in \mathbf{B}\) 且 \(b' \in \mathbf{B}'\)（相应地，若 \(c \in \mathbf{C}\) 且 \(c' \in \mathbf{C}'\)），F 中的元素 \(\langle b, b' \rangle\)（相应地，\(\langle c, c' \rangle\)）由 13.2.2 定义；于是有：

\[
\langle \iota a, b ^ {\prime} \rangle = \langle a, i b ^ {\prime} \rangle \quad \text { and } \quad \langle \sigma b, c ^ {\prime} \rangle = \langle b, s c ^ {\prime} \rangle \quad \text { if } \quad a \in \mathbf {A}, \quad b \in \mathbf {B}, \quad b ^ {\prime} \in \mathbf {B} ^ {\prime}, \quad c ^ {\prime} \in \mathbf {C} ^ {\prime}.
\]

13.3.5. 设 Y 是一个 \(C^{r}\) 类流形，\(\varphi: X \to Y\) 是一个 \(C^{r}\) 类态射，并设 \(x \in X\)，\(y \in Y\) 满足 \(y = \varphi(x)\)。13.2.3 中定义的映射 \(\varphi_{*}: T_{x}^{(r)}(X) \to T_{y}^{(r)}(Y)\) 与这些空间的滤链相容；通过转到相伴分次对象，它定义了一个线性映射

\[
\operatorname{gr} \left(\varphi_ {*}\right): \operatorname{gr} \mathrm{T} _ {x} ^ {(r)} (\mathrm{X}) \rightarrow \operatorname{gr} \mathrm{T} _ {y} ^ {(r)} (\mathrm{Y}).
\]

我们也用 \(\mathsf{TS}^{(r)}(\mathsf{T}_x(\varphi))\) 表示从 \(\mathsf{TS}^{(r)}(\mathsf{T}_x(\mathbf{X}))\) 到 \(\mathsf{TS}^{(r)}(\mathsf{T}_y(\mathbf{Y}))\) 的、由典范扩张 \(\mathsf{TS}(\mathsf{T}_x(\varphi))\)（\(\mathsf{T}_x(\varphi)\) 的典范扩张）诱导的映射。图表

\[
\begin{array}{c} \operatorname{gr} T _ {x} ^ {(r)} (X) \xrightarrow {\operatorname{gr} (\varphi_ {*})} \operatorname{gr} T _ {y} ^ {(r)} (Y) \\ \Biggl \downarrow i _ {\mathbf {X}, x} \\ \operatorname{TS} ^ {(r)} (T _ {x} (X)) \xrightarrow {\operatorname{TS} ^ {(r)} (T _ {x} (\varphi))} \operatorname{TS} ^ {(r)} (T _ {y} (Y)) \end{array}
\]

是交换的。

13.3.6. 假设 K 的特征为零。利用 A, IV, § 5, no. 8 中定义的同构 \(\varphi_{\mathbb{M}}: \mathsf{S}^k(\mathbf{M}) \to \mathsf{TS}^k(\mathbf{M})\)，可以在前述所有内容中把 \(\mathsf{TS}^k(\mathbf{T}_x(\mathbf{X}))\) 换成 \(k\) 次对称幂 \(\mathsf{S}^k(\mathbf{T}_x(\mathbf{X}))\)，即切空间 \(\mathbf{T}_x(\mathbf{X})\) 的对称幂。

